\par\noindent\textbf{(1)} 以界面z=0为参考，流体下上势扰动分别取φ1∝exp(kz)、φ2∝exp(−kz)，二者在远场衰减。运动学边界(∂t+Uj∂x)η=∂zφj，线化Bernoulli压力与表面Laplace压力差连接两侧幅度，要求频率共用。
\par\noindent\textbf{(2)} 消去振幅得到ρ1(ω−kU1)²+ρ2(ω−kU2)²=(ρ1−ρ2)gk+σk³。记ρΣ=ρ1+ρ2及质量加权流速Ubar=(ρ1U1+ρ2U2)/ρΣ，配方得(ω−kUbar)²=[(ρ1−ρ2)gk+σk³]/ρΣ−ρ1ρ2k²(U1−U2)²/ρΣ²。
\par\noindent\textbf{(3)} 根号内负时出现指数增长支。稳定要求|U1−U2|²≤[ρΣ/(ρ1ρ2)][(ρ1−ρ2)g/k+σk]。σ>0时右端在k*=√[(ρ1−ρ2)g/σ]最小，所有波数稳定的最大允许剪切平方为2ρΣ√[(ρ1−ρ2)gσ]/(ρ1ρ2)。
\par\noindent\textbf{(4)} 若σ=0，剪切项随k²而重力恢复随k，因此任何非零剪切在足够大波数下都可令判别式为负；理想无黏零厚度界面具有短波病态。真实流体有有限剪切层、黏性及表面张力，不能把任意短波不稳定直接当作真实材料预测。
